\section{What the game reads}\label{sec:discussion}

The game tests approximate invariance under adding a Bell pair, hence the
ratio two in the spectrum. This is a state-specific test. For a factor and its full commutant, so that Haag duality holds, the
classification uses the stronger statement that the value is one at
\emph{every faithful normal state}. Under the flow criterion of
Ref.~\cite[Section~4]{TG}, that statement is equivalent to
$1/2\in S(M_A)$: it holds for III$_1$ and also for
III$_\lambda$ with $\lambda=2^{-1/k}$ for a positive integer $k$.
A single fixed state with a very high score does not establish this
all-states property.

Using the ratio-three game and its corresponding factor criterion from Ref.~\cite{TG}, we can combine that game with $G$ by sending
both players a common tag telling them which test is being played.
With equal tag probabilities the resulting value is exactly
$(\omega_2+\omega_3)/2$, because strategies on the two tagged question
sets can be chosen independently. Value one at every faithful normal state
then requires both $1/2$ and $1/3$ in the factor invariant and singles out
III$_1$. The incommensurate-ratio mechanism is already part of the
flow-of-weights account of embezzlement~\cite{vLSWW24}; the tagged game is
an operational way of combining the tests.

Finite tensor powers show the same selectivity. For pairs with Schmidt
probabilities $(2/3,1/3)$, the Bell-pair error decreases, whereas ratios
three and four retain a positive defect (Table~\ref{tab:pairs}). Extensive
entropy is compatible with either behavior.

The positive floors follow directly from the multiplicative spacing.
If the nonzero probabilities lie in $c\,r^{\mathbb Z}$ with $r>2$, put
$a=1/\sqrt2$, $b=\sqrt{r/2}$ and
$t_m=(D\lambda)_m/\lambda_m$ wherever $\lambda_m>0$.
Each $t_m$ lies outside $(a,b)$, so
$(a+b)t_m\le t_m^2+ab$. Multiply by $\lambda_m^2$ and sum;
since $\sum_{\lambda_m>0}\lambda_m^2t_m^2\le1$, this gives
\[
\cE\ge2-2\frac{1+ab}{a+b}>0.
\]
For ratios three and four the bounds are respectively
$2-(1+\sqrt3)/\sqrt2=0.0681483\ldots$ and
$2-4\sqrt2/3=0.1143819\ldots$, uniformly in the number of copies.

\begin{table}[htb]
\centering\small
\begin{tabular}{@{}lcccccc@{}}
\toprule
Schmidt probabilities & Ratio & $n=1$ & $16$ & $64$ & $256$ & $1000$\\
\midrule
$(2/3,1/3)$ & 2 & 0.3905 & 0.1392 & 0.0417 & 0.0179 & 0.0091 \\
$(3/4,1/4)$ & 3 & 0.3270 & 0.1079 & 0.0718 & 0.0682 & 0.0681 \\
$(4/5,1/5)$ & 4 & 0.3029 & 0.1286 & 0.1146 & 0.1144 & 0.1144 \\
\bottomrule
\end{tabular}
\caption{Bell-pair embezzlement error for $n$ copies of the indicated
two-qubit state, computed with exact binomial multiplicities.}\label{tab:pairs}
\end{table}

The explicit strategy also clarifies the experimental requirements. In a
known Schmidt basis, the measurements reduce to reversible quotient and
remainder arithmetic followed by one fixed five-level measurement.
Their dependence on Schmidt rank is polynomial in its logarithm. Learning
or implementing that basis for a many-body state is a separate task;
entanglement-Hamiltonian spectroscopy and tomography address aspects of
it~\cite{Dalmonte18,Kokail21}. The finite sine-profile resource avoids
this step by choosing the computational basis from the outset.
Its separation from all flat-spectrum resources is quantitative and
survives the preparation error stated in Section~\ref{sec:operational-lower}.

The spectral counting criterion determines the sharp extraction cost and
its growing-target crossover. The original game has an assisted
certificate at $N=2^{19}$ with one shared $\Phi_9$, and a chain-only
certificate at $N=2^{84}$. The four-spin witness comes from tilted CHSH,
with classical play in the original game branch. The half-chain type
guarantees value one throughout the critical XX sector, and the modular
spectrum obstructs attainment at the ground state. The XXZ law concerns
the corner-transfer-matrix spectrum as the anisotropy approaches one
from above, with the physical spectral identification taken as input.
The cut-preserving bulk circuit retains
the finite critical XX Schmidt spectrum by exact conjugacy.

\section*{How this work was done}
The authors chose the questions and directed the work, which used AI
models from two developers: Claude Fable~5.2 (Anthropic) and GPT-6
Astra and GPT-6.1 Sol (OpenAI). Before release the manuscript was
checked refutation-first by fresh model instances that had not worked
on it: an earlier revision in three independent passes by Claude
Opus~5.5, two on the proofs and one on priority and credit, and the
following revision by GPT-6.1 Sol, which also reran the chain-only
certificate. This version applies the wording fixes from that check,
made with Claude Opus~5.5. The game constants, loss bounds and
separations rest on exact or outward-rounded certificates that the
accompanying archive re-checks; the other numerical values are
floating-point computations with the error control of
Appendix~\ref{app:numerics}. The theorems are proved in the text; none
is formalized in a proof assistant. The authors take responsibility for
every claim.

\section*{Reproducibility}
The accompanying source archive contains the numerical data, exact game certificates, and programs used to reproduce the calculations and figures. The game quadratic form is evaluated from all 243 question pairs, with its matrix identity checked in $\mathbb Q(\sqrt2)$. Numerical methods and truncation estimates are given in Appendix~\ref{app:numerics}.
